\documentclass[11pt,a4paper]{article}
\usepackage[margin=24mm]{geometry}
\usepackage{amsmath,amssymb,booktabs,array}
\usepackage[hidelinks]{hyperref}
\usepackage{microtype}
\newcommand{\dd}{\mathrm{d}}
\newcommand{\ii}{\mathrm{i}}
\newcommand{\cO}{\mathcal O}
\newcommand{\cG}{\mathcal G}
\title{Route G, Scheme A: One Internal Spectrum and Physical Probes}
\author{Research record: G0--G1}
\date{22 September 2026}
\begin{document}
\maketitle
\begin{abstract}
We construct a bounded, falsifiable realization of the common-spectrum
idea: a positive-energy complex scalar on a fixed circle, coupled to two
specified probes. One action determines both the mass tower and its
visibility. A passive basis change cannot change masses; a different
physical background can change them. Degenerate standing waves reveal
an important limitation: even two probes need not observe every state.
This is a conditional field-theory prototype, not a derivation of
elementary-particle masses, a string completion, or the full claims of
the motivating energy-spectrum document.
\end{abstract}

\section{Declared action, parameters and scope}
Use natural units and signature $(+----)$ on
$\mathbb M^{3,1}\times S^1$:
\begin{equation}
 \dd s^2=\eta_{\mu\nu}\dd x^\mu\dd x^\nu-R^2\dd\theta^2,\qquad
 \theta\sim\theta+2\pi,\quad R>0.
\end{equation}
The periodic complex scalar has an external flat connection
$A_\theta=\alpha$, $D_\theta=\partial_\theta+\ii\alpha$.
The gauge-invariant holonomy is $\exp(2\pi\ii\alpha)$; $\alpha$ is
defined modulo integers. The action is
\begin{equation}\label{eq:action}
 S_0=\int\dd^4x\int_0^{2\pi}R\,\dd\theta
 \left(\partial_\mu\Phi^\dagger\partial^\mu\Phi
 -\frac{|D_\theta\Phi|^2}{R^2}-M_5^2|\Phi|^2\right),
 \qquad M_5^2\geq0 .
\end{equation}
There is one common $R,M_5,\alpha$, not one parameter set per particle.
Their dimensions are $[R]=-1$, $[M_5]=1$, $[\alpha]=0$ and
$[\Phi]=3/2$. The connection and metric are prescribed backgrounds:
no dynamical gauge field, gravity, radion stabilization or preferred
holonomy is supplied. These omissions limit physical interpretation;
they do not invalidate a calculation conditional on \eqref{eq:action}.

This is an additional \emph{spatial} dimension of a specified model,
not a reinterpretation of a Spin(10) representation dimension.
The model does not identify its variables with the source document's
$E_i,E_p,E_n$, its symbol $R$, or negative Hamiltonian energies.
It retains a common-spectrum intuition from that document
\cite{chen}, using the conventional compactification mechanism
\cite{tong,yamada}. Route F is unchanged.

\section{Canonical reduction and positive energies}
The orthonormal eigenfunctions of the internal operator are
\begin{equation}
 u_n(\theta)=\frac{e^{\ii n\theta}}{\sqrt{2\pi R}},\qquad
 \int_0^{2\pi}R\,\dd\theta\,u_m^*u_n=\delta_{mn},\qquad
 -D_\theta^2u_n=(n+\alpha)^2u_n .
\end{equation}
Expanding $\Phi=\sum_{n\in\mathbb Z}\phi_n(x)u_n(\theta)$ in
\eqref{eq:action}, orthogonality gives the canonically normalized
four-dimensional action
\begin{equation}\label{eq:mass}
 S_0=\sum_n\int\dd^4x\left(
 |\partial_\mu\phi_n|^2-m_n^2|\phi_n|^2\right),\qquad
 \boxed{m_n^2=M_5^2+\frac{(n+\alpha)^2}{R^2}} .
\end{equation}
Here $[\phi_n]=1$. With sources switched off, the classical Hamiltonian is
\begin{equation}
 H=\sum_n\int\dd^3x\left(
 |\dot\phi_n|^2+|\nabla\phi_n|^2+m_n^2|\phi_n|^2\right)\geq0 .
\end{equation}
After quantization and the usual vacuum subtraction, particle and
antiparticle excitations have positive
$\omega_n(\boldsymbol p)=\sqrt{\boldsymbol p^2+m_n^2}$.
A negative Fourier label $n$ denotes compact momentum, not negative
energy. Opposite frequency solutions do not introduce negative-norm
states. The massless special case $M_5=0$, $n+\alpha=0$ requires the
usual care at zero external momentum.

For adjacent \emph{signed Fourier labels}, not adjacent sorted masses,
\begin{equation}\label{eq:diff}
 m_{n+1}^2-2m_n^2+m_{n-1}^2=\frac{2}{R^2},\qquad
 \Delta^3m_n^2=0 .
\end{equation}
This is a restrictive, testable tower relation. Four or more assigned
consecutive branches cannot have arbitrary independent masses.
A large light-to-heavy hierarchy is possible near a massless branch;
there is no claim that the circle forbids every hierarchy.

\section{Visibility derived from a fixed interaction}
Place probes at $0,\pi$ and choose the positive arc from $0$ to $\pi$.
Parallel transport to the first probe is
\begin{equation}
 W(0,\pi)=\exp\left(\ii\int_0^\pi A_\theta\dd\theta\right)
 =e^{\ii\pi\alpha},\qquad
 \cO_\pm=\frac{\Phi(0)\pm W(0,\pi)\Phi(\pi)}{\sqrt2}.
\end{equation}
Specify one coupling, with no mode-dependent projector choices:
\begin{equation}\label{eq:probe}
 S_{\rm probe}=\int\dd^4x\,
 \left[\lambda J_+^\dagger\cO_+
       +\lambda J_-^\dagger\cO_-+\mathrm{h.c.}\right].
\end{equation}
We take $[\lambda]=0$, $[J_\pm]=5/2$, $\lambda\ne0$.
The sources are external; this is linear response, not a dynamical
detector or a particle-conversion process.

Under $\Phi\mapsto e^{\ii\chi}\Phi$ and
$A_\theta\mapsto A_\theta-\partial_\theta\chi$,
\begin{equation}
 W\mapsto e^{-\ii(\chi(\pi)-\chi(0))}W,\qquad
 \cO_\pm\mapsto e^{\ii\chi(0)}\cO_\pm .
\end{equation}
The operators are \emph{charged and gauge-covariant}, not local
gauge-invariant observables of a dynamical gauge theory. Transforming
$J_\pm$ with the same charge makes \eqref{eq:probe} invariant.
Changing the chosen path in a nontrivial holonomy changes the physical
probe definition, not merely a gauge convention.

Writing $\lambda\cO_a=\sum_n g_{an}\phi_n$ gives
\begin{align}
 g_{\pm n}&=\frac{\lambda}{\sqrt{4\pi R}}
             (1\pm e^{\ii\pi(n+\alpha)}),\\
 \rho^{(n)}_{ab}&=g_{an}g_{bn}^*,\qquad \rho^{(n)}\succeq0,\\
 |g_{\pm n}|^2&=\frac{|\lambda|^2}{2\pi R}
                  \left[1\pm\cos\pi(n+\alpha)\right],\\
 \rho^{(n)}_{+-}&=\frac{\ii|\lambda|^2}{2\pi R}
                   \sin\pi(n+\alpha),\\
 \operatorname{tr}\rho^{(n)}&=\frac{|\lambda|^2}{\pi R}.
 \label{eq:sum}
\end{align}
The normalized weights
$p_{\pm n}=(1\pm\cos\pi(n+\alpha))/2$ sum to one, but are
\emph{relative probe strengths}, not branching ratios.

For $Q^2\geq0$, away from a massless zero-momentum singularity,
the Euclidean source response is
\begin{equation}\label{eq:response}
 \cG_{ab}(Q^2)=\sum_{n\in\mathbb Z}
 \frac{\rho^{(n)}_{ab}}{Q^2+m_n^2}.
\end{equation}
Its continuation has the same possible poles in every channel;
a zero residue can hide a pole from a particular channel.
For every source vector $v$,
$v^\dagger\cG v=\sum_n|g_n^\dagger v|^2/(Q^2+m_n^2)\geq0$.
The response has dimension $-1$ in these conventions.

\subsection*{Degeneracy: a genuine dark subspace}
For a mass eigenvalue $s$, let $C_s$ have columns $g_n$ for all
$n$ with $m_n^2=s$. The invariant pole residue is
\begin{equation}\label{eq:deg}
 R_s=C_sC_s^\dagger,\qquad
 \dim\ker C_s=d_s-\operatorname{rank}C_s,\qquad
 \operatorname{tr}R_s=d_s\frac{|\lambda|^2}{\pi R}>0 .
\end{equation}
The branch sum \eqref{eq:sum} holds in the Fourier basis; it is
\emph{not} a claim of equal strength for every superposition in a
degenerate eigenspace. At $\alpha=0$, the $n=k,-k$ columns coincide
for $k\ne0$. Their antisymmetric standing wave
\begin{equation}
 \frac{u_k-u_{-k}}{\sqrt2}
  =\frac{\ii\sin(k\theta)}{\sqrt{\pi R}}
\end{equation}
vanishes at both probe positions. Thus $\operatorname{rank}C_s=1$
and a massive state can be dark to \emph{both} probes.
The orthogonal combination remains visible: the entire mass pole
does not disappear from both channels. At $\alpha=1/2$, the
degenerate pair $n,-n-1$ has independent (indeed orthogonal) coupling
columns, so this two-dimensional dark subspace obstruction is absent.
This is an apparatus-dependent result, not a universal rule of nature.

\section{What does and does not change a mass}
In any retained mode space, let $K=Q^2I+\mathsf M^2$ and let $B$ have
columns $g_n$. Under $\phi=U\widetilde\phi$ with $U$ unitary,
\begin{equation}
 \widetilde K=U^\dagger K U,\qquad
 \widetilde B=BU,\qquad
 \widetilde B\,\widetilde K^{-1}\widetilde B^\dagger
 =BK^{-1}B^\dagger .
\end{equation}
The masses and the transformed response do not change. Rotating
only a mass matrix or only a probe is not a passive change of the
same experiment. For a general invertible field redefinition the
kinetic metric must also be transformed.

An integer large gauge transformation is
\begin{equation}
 \alpha\mapsto\alpha+\ell,\quad
 \Phi\mapsto e^{-\ii\ell\theta}\Phi,\quad
 n\mapsto n-\ell .
\end{equation}
Both $n+\alpha$ and the probe coefficients remain unchanged.
A finite cutoff comparison must relabel the retained window too.
By contrast, physically changing $R$ or a gauge-inequivalent holonomy
can change the spectrum. Not every distinct value of $\alpha$ gives
a distinct mass set: $\alpha$ and $-\alpha$ have the same mass set
after reflection of labels. A comparison of two static backgrounds
does not describe the energy cost, vacuum selection, or time
evolution of such a change.

\section{Numerical verification and truncation control}
No experimental masses are input. The four engineering cards have
$M_5=0.5$ and $\lambda=1$ in common arbitrary energy units.
For the branch $n=0$:
\begin{center}
\begin{tabular}{ccccc}
\toprule
$R$ & $\alpha$ & $m_0^2$ & $p_{+0}$ & $p_{-0}$\\
\midrule
1 & 0 & 0.250000 & 1 & 0\\
1 & 0.25 & 0.312500 & 0.853553 & 0.146447\\
1 & 0.5 & 0.500000 & 0.5 & 0.5\\
2 & 0.25 & 0.265625 & 0.853553 & 0.146447\\
\bottomrule
\end{tabular}
\end{center}
The decimals display exact formulae, not fitted coefficients.
At $\alpha=0,n=1$, $m_1^2=1.25$ but $p_{+1}=0$:
mass and visibility are different observables.

An independent matrix discretization uses $h=2\pi/N$,
$\theta_j=jh$, $U_j=e^{\ii\alpha h}$ and periodic indices:
\begin{equation}
 (L_h\psi)_j=M_5^2\psi_j+
 \frac{2\psi_j-U_j\psi_{j+1}-U_{j-1}^*\psi_{j-1}}{R^2h^2}.
\end{equation}
Its quadratic form is
$Rh\sum_j[M_5^2|\psi_j|^2+
|U_j\psi_{j+1}-\psi_j|^2/(R^2h^2)]\geq0$.
The matrix eigenvalues are checked against
\begin{equation}
 m_{n,h}^2=M_5^2+
 \frac{4}{R^2h^2}\sin^2\left(\frac{(n+\alpha)h}{2}\right)
 =m_n^2-\frac{(n+\alpha)^4h^2}{12R^2}+O(h^4).
\end{equation}
Only resolved low labels are used for continuum convergence;
the lattice's finite Brillouin zone is not an extra prediction.
For a small gauge change, $U_j\mapsto
e^{\ii\chi_j}U_je^{-\ii\chi_{j+1}}$, transform both the matrix
and the discrete Wilson-dressed probes. Degenerate comparisons use
\eqref{eq:deg}, not individual numerical eigenvectors.

For the symmetric mode truncation $|n|\leq N_{\rm mode}>|\alpha|$,
every response matrix entry obeys the conservative bound
\begin{align}
 |\cG_{ab}-\cG^{(N_{\rm mode})}_{ab}|
 &\leq \frac{|\lambda|^2 R}{\pi}
       \sum_{|n|>N_{\rm mode}}\frac{1}{(|n|-|\alpha|)^2}\\
 &\leq\boxed{\frac{2|\lambda|^2R}
                   {\pi(N_{\rm mode}-|\alpha|)}} .
\end{align}
This follows from $|\rho^{(n)}_{ab}|\leq|\lambda|^2/(\pi R)$,
$Q^2,M_5^2\geq0$ and the integral bound for a decreasing positive
series. A difference between two finite cutoffs is not itself the
exact infinite-tail error. The implementation records both successive
differences and this a priori bound.

An independent infinite-sum check follows from Poisson summation.
For $\beta=R\sqrt{M_5^2+Q^2}>0$, Fourier transforming
$(q^2+\beta^2)^{-1}$ yields
\begin{equation}
 F(x)=\sum_n\frac{e^{\ii(n+\alpha)x}}{(n+\alpha)^2+\beta^2}
 =\frac{\pi}{\beta}\sum_{\ell\in\mathbb Z}
 e^{2\pi\ii\ell\alpha}e^{-\beta|x-2\pi\ell|}.
\end{equation}
Set $a=e^{-\pi\beta}$, $t=2\pi\alpha$ and
$D=1+a^4-2a^2\cos t$. Geometric summation at $x=0,\pi$ gives
\begin{equation}
 S_0=F(0)=\frac{\pi}{\beta}\frac{1-a^4}{D},\qquad
 S_\pi=F(\pi)=\frac{\pi}{\beta}
       \frac{a(1-a^2)(1+e^{\ii t})}{D},
\end{equation}
and therefore
\begin{equation}
 \cG=\frac{|\lambda|^2R}{2\pi}
 \begin{pmatrix}
 S_0+\operatorname{Re}S_\pi & \ii\operatorname{Im}S_\pi\\
 -\ii\operatorname{Im}S_\pi & S_0-\operatorname{Re}S_\pi
 \end{pmatrix}.
\end{equation}
All four cards have $\beta>0$ at the tested momenta. This formula
provides an exact response reference, not a large-cutoff surrogate.

The executable and machine-readable report are respectively
\begin{center}\small
\texttt{route\_g/code/verify\_g1\_circle\_spectrum.py}\\
\texttt{route\_g/output/g1\_circle\_spectrum.json}.
\end{center}
They record matrix convergence, full residues, gauge/basis covariance,
dark-subspace ranks, positivity, truncations and source hash.
Checks establish implementation consistency, not experimental evidence.

\section{Next physical gate: an actual transition}
The free action is diagonal in $n$. It supplies no spontaneous
$\phi_n\to\phi_m$ conversion. A prescribed external source may inject
energy and excite a mode, but is not yet an energy-accounted dynamical
detector.

Even a natural bulk extension can fail before any lengthy calculation.
Consider a massless periodic neutral scalar $\chi$ and the candidate
vertex $-\lambda_5\chi|\Phi|^2$ (not a completed stable potential).
Its Fourier overlap enforces
$n=m+q$, and $m_{\chi_q}=|q|/R$. For $M_5>0$ and $n\ne m$,
the strict Euclidean triangle inequality for
$(M_5,(m+\alpha)/R)$ and $(0,(n-m)/R)$ gives
\begin{equation}\label{eq:closed}
 \sqrt{M_5^2+\frac{(n+\alpha)^2}{R^2}}
 <
 \sqrt{M_5^2+\frac{(m+\alpha)^2}{R^2}}+
 \frac{|n-m|}{R}.
\end{equation}
Hence $\phi_n\to\phi_m+\chi_q$ is kinematically closed.
For $M_5=0$ the non-strict inequality still holds; equality supplies
only threshold kinematics, not an ordinary open two-body decay
phase space. Other processes, scattering, extra fields or modified
backgrounds are not excluded by this particular argument.

G2 should therefore choose \emph{one} localized dynamical
interaction/detector, or one explicitly energy-accounted driven
background, and compute one allowed process. Localization is a new
physical assumption: compact momentum must be exchanged with that
structure, not silently violated. G3 may study background selection;
G4 must derive spin, charges and chirality before a particle
identification or a Route-F bridge is proposed.

\paragraph{Bounded conclusion.}
G1 realizes one common positive spectrum with interaction-defined
visibility and exact correlated tower relations. It rejects
coordinate-induced mass conversion, while allowing genuine dark
standing waves for an incomplete apparatus. It does not repair
Route-F flavor/seesaw tensions, derive three chiral families, or prove
the motivating document. The next useful step is a transition
amplitude with a stated source of energy and momentum, not another
mass fit.

\begin{thebibliography}{9}
\bibitem{chen}
B.~Chen, \emph{Energy spectrum for elementary particles}, final
conference document (2022), five pages.
\href{https://indico.cern.ch/event/1109513/contributions/4969967/attachments/2480973/4410560/Energy%20spectrum%20for%20elementary%20particles_final.pdf}{Exact final PDF}.
Source checksum and interpretation ledger: \texttt{route\_g/SOURCES.md}.
\bibitem{tong}
D.~Tong, \emph{Lectures on String Theory}, section 8,
\href{https://arxiv.org/html/0908.0333v3#S8}{arXiv:0908.0333}.
No string oscillator or winding sector is assumed here.
\bibitem{yamada}
A.~Yamada, \emph{The UV sensitivity of the Higgs potential in
Gauge--Higgs Unification}, PTEP 2021, 093B01, section 2, equation (10),
\href{https://doi.org/10.1093/ptep/ptab085}{doi:10.1093/ptep/ptab085}.
\end{thebibliography}
\end{document}
